\documentclass[11pt,a4paper]{article}
\usepackage{fontspec}
\setmainfont{Latin Modern Roman}
\usepackage[margin=23mm]{geometry}
\usepackage{amsmath,amssymb,amsthm,mathtools,mathrsfs}
\usepackage{longtable,booktabs}
\usepackage[unicode,hidelinks]{hyperref}
\providecommand{\C}{\mathbb C}
\providecommand{\R}{\mathbb R}
\newtheorem{theorem}{Theorem}[section]
\newtheorem{lemma}[theorem]{Lemma}
\newtheorem{proposition}[theorem]{Proposition}
\newtheorem{corollary}[theorem]{Corollary}
\newtheorem{definition}[theorem]{Definition}
\newtheorem{remark}[theorem]{Remark}
\allowdisplaybreaks
\emergencystretch=3em
\makeatletter
\renewcommand\@pnumwidth{2.5em}
\renewcommand\@tocrmarg{3.5em}
\makeatother
\title{The Original Arithmetic and Gamma Mixed Transfer}
\author{Split-Zero research programme}
\date{14 September 2026}
\begin{document}
\maketitle
\paragraph{Portable source edition.}
Every TeX input used by this entry is expanded below. Historical file
and directory names in the preserved text are source locators; the
exact origin paths and hashes are recorded in \texttt{11\_SOURCE\_MAP.json}.
The other numbered proof files in this flat handoff supply the
accompanying complete calculations.\par

% Exact original arithmetic mixed transfer. Historical sources remain unchanged.
\section{The original arithmetic kernel and boundary receive the full source comparison}
\label{sec:AMT}

This calculation uses the proved full-polynomial comparison TW12/TWA3,
the original observation OPG1--5, and the complete original Gamma
kernel estimate OKV1--27. It gives the exact maps between their source,
quotient, kernel and boundary metrics, before taking any determinant.
The subsequent simple-quartet conclusion uses the evaluated original
period test RPC1--26 and its exact rank calculation OQX1--3. Every
source order below is one of the already present orders $1,k$.

\subsection{The complete original source and the same observation}
Retain the original full-order quartet, with $0<\delta<1/2$,
$\gamma>2$, and positive integer complete zero order $m$. Thus
\[
\begin{gathered}
 \mathcal R=\{\tfrac12+\delta+i\gamma,\tfrac12+\delta-i\gamma,
              \tfrac12-\delta+i\gamma,\tfrac12-\delta-i\gamma\},\quad
 h(z)=\prod_{\alpha\in\mathcal R}(z-\alpha)^m,\quad v_h=2\xi/h,\\
 k\ge3,\quad c=k/2,\quad e=1+k(m-1),\quad q=e(k+1)^2,\\
 \lambda_{a,b}=c+(2a-k)\delta+i(2b-k)\gamma,\quad
 \chi(S)=\prod_{a,b=0}^k(S-\lambda_{a,b})^e,\quad E=\mathbb C[S]/(\chi),\\
 w_h(y)=|v_h(1/2+iy)|^2/(2\pi),\quad m_{h,k}=w_h^{*k},\quad
 d\sigma(y)=|\Gamma(1/4+iy/2)|^2dy/(2\pi).
\end{gathered}\tag{AMT1}
\]
The original source inner products, conjugate-linear in their first
argument, are
\[
 \langle P,Q\rangle_{\rm ar}
  =\int_{\mathbb R}\overline{P(c+iy)}Q(c+iy)m_{h,k}(y)\,dy,
 \qquad
 \langle P,Q\rangle_\sigma
  =\int_{\mathbb R}\overline{P(c+iy)}Q(c+iy)\,d\sigma(y).
 \tag{AMT2}
\]
No measure mass or complex-coordinate phase is changed. In particular
$\int d\sigma=\sqrt{2\pi}$. Write $\mathcal P_N$ for polynomials
in the original $S$ of degree at most $N$. In its coefficient frame
$(1,S,\ldots,S^N)$ let the two Grams be $H_N^{\rm ar},H_N^\sigma$.
For $N\ge q-1$, the literal remainder map
$J_N:\mathcal P_N\to E$ is onto and has kernel
$\chi\mathcal P_{N-q}$; at $N=q-1$ this is the zero vector space.
The frame on $E$ is $(1,S,\ldots,S^{q-1})$.

Here is the original observation being retained. Put $n=4m$,
$E_h=\mathbb C[z]/(h)$, $U=M_{j_hv_h}$, and
$\iota[P]=[P(z_1+\cdots+z_k)]\in E_h^{\otimes k}$.
The map $j_h$ records every divided Taylor coefficient through order
$m-1$ at every original root. Its unit $j_hv_h$ is invertible because
the stated order is complete. Fix the original $u\ne0$ and its
logarithm branch. With $\Phi_h'=h$ and $\Phi_h(0)=0$, let
$\ell_j$ be the outward ray of angle
$(\arg u+\pi+2\pi j)/(n+1)$ and $\Gamma_j=-\ell_0+\ell_j$.
The period map and projector are exactly
\[
\begin{gathered}
 (\Pi[P])_j=\int_{\Gamma_j}e^{\Phi_h(z)/u}
                           \operatorname{rem}_hP(z)\,dz,\quad1\le j\le n,\\
 C:\Gamma_j\mapsto\Gamma_{j+1}-\Gamma_1\ (j<n),\quad
        \Gamma_n\mapsto-\Gamma_1,\qquad
 P_k=\frac1{n+1}\sum_{a=0}^{n}(C^{-a})^{\otimes k},\\
 L=P_k\Pi^{\otimes k}U^{\otimes k}\iota,\quad
 B=\operatorname{im}L,\quad\jmath:B\hookrightarrow(\mathbb C^n)^{\otimes k},
 \quad\jmath\Lambda=L,\quad\Lambda:E\twoheadrightarrow B,\quad K=\ker\Lambda.
\end{gathered}\tag{AMT3}
\]
These are the convergent original representative integrals; their
existence and full determinant are proved in PDE and OPG. Fix the
original coefficient inclusion $I:K\hookrightarrow E$ and section
$S_*:B\to E$ with $\Lambda S_*=I_B$. Let
$r=\dim K$ and $b=\dim B=q-r$. All matrices below use these same
frames and this same full unit and period observation. Empty Grams
have determinant one, and maps to or from a zero space have their
unique empty values.

\subsection{All constants of the actual full-polynomial inequality}
Retain exactly the constants of TW and TWA:
\[
\begin{gathered}
 \alpha=\pi/2,\quad p=8m-\tfrac92,\quad B_h=42+8m,\quad M=21+4m,\\
 \vartheta_h=\int_{-1}^1w_h(y)\,dy>0,\quad
 M_\alpha=\int_{\mathbb R}e^{\alpha y}w_h(y)\,dy\in(0,\infty),\\
 c_\Gamma=\sqrt2e^{-7/6},\quad C_\Gamma=\sqrt{2\sqrt5}e^{2/3},\\
 a_k=\frac{c_h\vartheta_h^{k-3}e^{-\alpha(k-3)}(k-2)^{-B_h}}
                   {C_\Gamma2^{B_h/2}},\qquad
 A_k=\frac{C_h^{\rm tilt}}{c_\Gamma}k^{p+1}M_\alpha^{k-1},\qquad
 D_N=[1+4(N+M)^2]^M.
\end{gathered}\tag{AMT4}
\]
Here $c_h$ is the proved original convolution lower-envelope constant
TW7 and $C_h^{\rm tilt}$ is the original compact/tail maximum BT12.
Their full constructions occur in the preserved arithmetic endpoint
and BT providers; they depend on this fixed $h$ and not on $k,N,u$.
Their defining inequalities, proved there with the original numerator,
are
\[
\begin{split}
 m_{h,k}(y)&\ge c_h\vartheta_h^{k-3}
    e^{-\alpha(|y|+k-3)}(k-2+|y|)^{-B_h},\\
 m_{h,k}(y)&\le k C_h^{\rm tilt}M_\alpha^{k-1}
              e^{-\alpha|y|}(1+|y|/k)^{-p},\\
 c_\Gamma e^{-\alpha|y|}(1+|y|)^{-1/2}
 &\le d\sigma/dy\le
 C_\Gamma e^{-\alpha|y|}(1+|y|)^{-1/2}.
\end{split}\tag{AMT5}
\]
For clarity the full polynomial step is given here. The first and
third bounds, $k-2+|y|\le(k-2)(1+|y|)$, and
$(1+|y|)^{B_h}\le2^{B_h/2}(1+y^2)^M$, give
$m_{h,k}\ge a_k(1+y^2)^{-M}d\sigma/dy$. The second and third,
using $1+|y|/k\ge(1+|y|)/k$ and $p>1/2$, give
$m_{h,k}\le A_kd\sigma/dy$.
The fixed Gamma orthonormal polynomials obey
\[
 ye_j=\sqrt{(j+1)(j+1/2)}e_{j+1}
                         +\sqrt{j(j-1/2)}e_{j-1}.
\]
The two shifts on degree at most $N$ have norms at most $N+1$ and
$N$, respectively. Thus $\|yP\|_\sigma\le2(N+1)\|P\|_\sigma$.
Iterating at the successively increased degrees and expanding
$(1+y^2)^M$ proves
$\int|P|^2(1+y^2)^M d\sigma\le D_N\int|P|^2d\sigma$.
Cauchy--Schwarz applied to $|P|(1+y^2)^{-M/2}$ and
$|P|(1+y^2)^{M/2}$ then gives the reciprocal lower bound. The
coordinate map $P(S)\mapsto P(c+iy)$ has inverse
$f(y)\mapsto f((S-c)/i)$, preserves the degree and retains every
coefficient through this explicitly invertible substitution. It
therefore proves on the complete original source
\[
 \boxed{\quad \frac{a_k}{D_N}H_N^\sigma
                  \preceq H_N^{\rm ar}\preceq A_kH_N^\sigma.\quad}
 \tag{AMT6}
\]
All coefficients can be complex in this proof. In particular AMT6
holds on every affine remainder fibre, not only on monic vectors.
The two Grams are positive: the Gamma density is positive and the
lower bound in AMT5 is positive everywhere. Their finite moments
exist by the exponential upper bound.

\subsection{Exact remainder, minimum-section, kernel and boundary maps}
For $a\in\{\sigma,\mathrm{ar}\}$ put
\[
\begin{gathered}
 G_N^a=(J_N(H_N^a)^{-1}J_N^*)^{-1},\qquad
 V_N^a=(H_N^a)^{-1}J_N^*G_N^a:E\longrightarrow\mathcal P_N,\\
 H_{K,N}^a=I^*G_N^aI,\qquad
 Q_N^a=(\Lambda(G_N^a)^{-1}\Lambda^*)^{-1},\qquad
 S_N^a=(G_N^a)^{-1}\Lambda^*Q_N^a:B\longrightarrow E.
\end{gathered}\tag{AMT7}
\]
Inverses are on the indicated nonzero spaces. For an onto matrix
$J$ and positive $H$, direct multiplication gives $JV=I$ for
$V=H^{-1}J^*(JH^{-1}J^*)^{-1}$. Moreover
$Z^*HV=Z^*J^*G=0$ for $JZ=0$. Therefore every lift of $z$ is
$Vz+Z$, with squared norm
$z^*Gz+Z^*HZ$, proving its unique minimum and its Gram. Applying
this proof first to $J_N$, then to $\Lambda$, proves every formula
AMT7 and $\Lambda S_N^a=I_B$, $I^*G_N^aS_N^a=0$.
It also proves that $V_N^aS_N^a$ is the minimum lift for the
composite $\Lambda J_N$.

The exact comparison map on $\mathcal P_N$ is the identity on
original coefficients; it commutes with $J_N$ and $\Lambda J_N$.
It induces the identity on $E,K,B$. Between the two spaces of
minimum lifts the exact bijection and its inverse are
\[
 \operatorname{im}V_N^\sigma\longrightarrow\operatorname{im}V_N^{\rm ar},
 \quad V_N^\sigma z\longmapsto V_N^{\rm ar}z,
 \qquad T_N=V_N^{\rm ar}J_N\big|_{\operatorname{im}V_N^\sigma},
 \quad T_N^{-1}=V_N^\sigma J_N\big|_{\operatorname{im}V_N^{\rm ar}}.
 \tag{AMT8}
\]
Their difference $V_N^{\rm ar}z-V_N^\sigma z$ lies in
$\chi\mathcal P_{N-q}$, with its unique coefficient obtained by
division by the unchanged monic $\chi$. This also gives the
minimum-lift bijection on the preimage of $K$ and the analogous
bijection $V_N^\sigma S_N^\sigma y\mapsto
V_N^{\rm ar}S_N^{\rm ar}y$ for boundary lifts. The latter
differences lie in $\ker(\Lambda J_N)$, exactly as multiplication
by $\Lambda J_N$ verifies.

Minimize both sides of AMT6 over $J_NP=z$. Taking infima preserves
each inequality because these are the same nonempty fibres. Restrict
the resulting inequality to $z=Ix$, and independently minimize it
over $\Lambda z=y$. This proves the three full-vector comparisons
\[
\boxed{\begin{gathered}
 (a_k/D_N)G_N^\sigma\preceq G_N^{\rm ar}\preceq A_kG_N^\sigma,\\
 (a_k/D_N)H_{K,N}^\sigma\preceq H_{K,N}^{\rm ar}
                                      \preceq A_kH_{K,N}^\sigma,\\
 (a_k/D_N)Q_N^\sigma\preceq Q_N^{\rm ar}\preceq A_kQ_N^\sigma.
\end{gathered}}\tag{AMT9}
\]
The source identity is an exact bounded bijection with the two
constants in AMT6. No invariance of $K$ under an additional
multiplier was used to obtain its induced identities.

The matrix $[I,S_*]$ is invertible: its inverse sends $z$ to the
unique kernel coordinate of $z-S_*\Lambda z$ and the boundary
coordinate $\Lambda z$. Each $S_N^a-S_*$ lies in $K$, so the
change from this frame to $[I,S_N^a]$ is triangular with identity
diagonal and determinant one. The latter frame is orthogonal in
$G_N^a$ by AMT7. Consequently
\[
 \boxed{\quad\det H_{K,N}^a\det Q_N^a
       =|\det[I,S_*]|^2\det G_N^a,\qquad a\in\{\sigma,\mathrm{ar}\}.\quad}
 \tag{AMT10}
\]
This includes $r=0$ and $b=0$ and retains the original frame factor.

\subsection{The exact nested deficits and their rank factors}
Define the actual finite source determinant deficit and its width
\[
 \Xi_{k,N}=(N+1)\log A_k-
                  \log\frac{\det H_N^{\rm ar}}{\det H_N^\sigma},
 \qquad L_N=\log(A_kD_N/a_k)\ge0.
 \tag{AMT11}
\]
The last sign follows by applying AMT6 to any nonzero vector.
Put
\[
\begin{split}
 d_{E,N}&=q\log A_k-\log\frac{\det G_N^{\rm ar}}{\det G_N^\sigma},\\
 d_{K,N}&=r\log A_k-\log\frac{\det H_{K,N}^{\rm ar}}{\det H_{K,N}^\sigma},\\
 d_{B,N}&=b\log A_k-\log\frac{\det Q_N^{\rm ar}}{\det Q_N^\sigma}.
\end{split}\tag{AMT12}
\]
AMT9, by congruence with the inverse square root of each comparison
Gram, places each eigenvalue of its relative Gram in
$[a_k/D_N,A_k]$. Multiplication of all eigenvalues gives
$0\le d_{K,N}\le rL_N$, $0\le d_{B,N}\le bL_N$ and
$0\le d_{E,N}\le qL_N$.

Here is the full source-to-quotient determinant calculation, including
all relation directions. Choose an $H_N^\sigma$-orthonormal
coefficient frame $R_N$ of $\ker J_N$ solely for this computation,
and put $W_N=[R_N,V_N^\sigma(G_N^\sigma)^{-1/2}]$.
AMT7 gives $W_N^*H_N^\sigma W_N=I_{N+1}$. In this frame the
positive contraction
\[
 \mathscr D_N=A_k^{-1}W_N^*H_N^{\rm ar}W_N
       =\begin{pmatrix}F_N&C_N'\\(C_N')^*&A_N'\end{pmatrix}
       \preceq I_{N+1}
\]
has Schur complement
\[
 T_N'=A_N'-(C_N')^*F_N^{-1}C_N'
       =A_k^{-1}(G_N^\sigma)^{-1/2}G_N^{\rm ar}(G_N^\sigma)^{-1/2}.
\]
Indeed minimization in its first block is exactly the original
remainder minimum at $z=(G_N^\sigma)^{-1/2}w$.
Both $F_N$ and $T_N'$ are positive contractions: positivity follows
from positive block elimination, and the upper bound for the Schur
complement follows by testing its minimum at zero first coordinate.
Since $|\det W_N|^2=(\det H_N^\sigma)^{-1}$, elimination gives
\[
 \Xi_{k,N}=d_{R,N}+d_{E,N},\quad d_{R,N}:=-\log\det F_N\ge0,
 \qquad d_{E,N}=-\log\det T_N'\ge0.
 \tag{AMT13}
\]
When $N=q-1$, $F_N$ is empty, $d_{R,N}=0$, and the same equalities
hold without a relation inverse. The determinant identities are
independent of the auxiliary orthonormal frame.
Dividing the two instances of AMT10 proves
$d_{K,N}+d_{B,N}=d_{E,N}$. Thus the complete nested statement is
\[
\boxed{\begin{gathered}
 \Xi_{k,N}=d_{R,N}+d_{K,N}+d_{B,N},\qquad d_{R,N},d_{K,N},d_{B,N}\ge0,\\
 0\le d_{K,N}\le\min\{rL_N,d_{E,N}\}\le\min\{rL_N,\Xi_{k,N}\},\\
 0\le d_{B,N}\le\min\{bL_N,d_{E,N}\}\le\min\{bL_N,\Xi_{k,N}\}.
\end{gathered}}\tag{AMT14}
\]
In particular each metric keeps its own rank factor, and the two
deficits retain their exact common sum. Any proved finite upper
bound for the original $\Xi_{k,N}$, including the degree-specific
HMR44/TWA15 bound, can be substituted in these displayed minima.
No new density or source determinant is introduced by this substitution.

\subsection{All four original endpoints and the correlated mixed return}
Let
\[
 (N_0,N_1,N_2,N_3)=(q-1,q,2q-1,2q),\qquad
 (\epsilon_0,\epsilon_1,\epsilon_2,\epsilon_3)=(1,1,-1,-1),
\]
and define, in their original signs,
\[
\begin{split}
 F_K^a&=\sum_{j=0}^3\epsilon_j\log\det H_{K,N_j}^a,\qquad
 F_B^a=\sum_{j=0}^3\epsilon_j\log\det Q_{N_j}^a,\\
 \mathcal B_k^a&=\sum_{j=0}^3\epsilon_j\log\det G_{N_j}^a,
 \quad\Delta_K^\sigma=F_K^{\rm ar}-F_K^\sigma,
 \quad\Delta_B^\sigma=F_B^{\rm ar}-F_B^\sigma,
 \quad\delta_k^\sigma=\mathcal B_k^{\rm ar}-\mathcal B_k^\sigma.
\end{split}\tag{AMT15}
\]
AMT10 cancels its four identical frame factors only after these
signs are applied. It gives
$F_K^a+F_B^a=\mathcal B_k^a$ and
$\Delta_K^\sigma+\Delta_B^\sigma=\delta_k^\sigma$ exactly.
The spaces $\mathcal P_N$ increase with the same source norm and
remainder map, so their quotient minima decrease. Restriction to
$K$ and minimization onto $B$ preserve that order. Pairing $N_0$
with $N_2$ and $N_1$ with $N_3$ consequently proves
$F_K^a,F_B^a\ge0$.

For $X=K,B$, put $r_K=r$, $r_B=b$, and retain the exact caps
\[
 c_{X,N}=\min\{r_XL_N,d_{E,N}\},\qquad
 C_X^-=c_{X,q-1}+c_{X,q},\qquad
 C_X^+=c_{X,2q-1}+c_{X,2q}.
 \tag{AMT16}
\]
The entirely source-evaluated caps
$\widehat c_{X,N}=\min\{r_XL_N,\Xi_{k,N}\}$ and the corresponding
$\widehat C_X^\pm$ are at least these values; all conclusions below
also hold with the hats. Substituting AMT12 in AMT15 records each
endpoint and its sign:
\[
\boxed{\begin{split}
 \Delta_K^\sigma&=-d_{K,q-1}-d_{K,q}+d_{K,2q-1}+d_{K,2q},\\
 \Delta_B^\sigma&=-d_{B,q-1}-d_{B,q}+d_{B,2q-1}+d_{B,2q},\\
 \delta_k^\sigma&=-d_{E,q-1}-d_{E,q}+d_{E,2q-1}+d_{E,2q},\\
 -C_X^-&\le\Delta_X^\sigma\le C_X^+,\qquad X=K,B.
\end{split}}\tag{AMT17}
\]
The four copies of $r_X\log A_k$ cancel here because their signs
sum to zero. Since the original total correction is retained, the
joint interval is stronger than two independent intervals:
\[
\boxed{\begin{split}
 \max\{-C_K^-,\delta_k^\sigma-C_B^+\}
   &\le\Delta_K^\sigma\le
          \min\{C_K^+,\delta_k^\sigma+C_B^-\},\\
 \max\{-C_B^-,\delta_k^\sigma-C_K^+\}
   &\le\Delta_B^\sigma\le
          \min\{C_B^+,\delta_k^\sigma+C_K^-\}.
\end{split}}\tag{AMT18}
\]
These follow by subtracting the interval for either component from
the exact common sum; they require no independence assertion.

Here are the explicit original four endpoint widths:
\[
\begin{split}
 E_k^+&=L_{q-1}+L_q
   =2\log(A_k/a_k)+M\log[1+4(q-1+M)^2]
                         +M\log[1+4(q+M)^2],\\
 E_k^-&=L_{2q-1}+L_{2q}
   =2\log(A_k/a_k)+M\log[1+4(2q-1+M)^2]
                         +M\log[1+4(2q+M)^2].
\end{split}\tag{AMT19}
\]
They are the unchanged TWA5 constants, both nonnegative. AMT17
therefore gives in particular
\[
 \boxed{-rE_k^+\le F_K^{\rm ar}-F_K^\sigma\le rE_k^-,\qquad
        -bE_k^+\le F_B^{\rm ar}-F_B^\sigma\le bE_k^-.}
 \tag{AMT20}
\]
The total correction retains TWA15 with all four individual source
deficits:
\[
\begin{split}
 \underline\delta_k^\sigma
   &=\max\{-qE_k^+,-\Xi_{k,q-1}-\Xi_{k,q}\},\\
 \overline\delta_k^\sigma
   &=\min\{qE_k^-,\Xi_{k,2q-1}+\Xi_{k,2q}\},\qquad
 \underline\delta_k^\sigma\le\delta_k^\sigma\le\overline\delta_k^\sigma.
\end{split}\tag{AMT21}
\]
This also follows immediately from AMT14 and AMT17. It can be
inserted into AMT18 when an interval, rather than the actual finite
total determinant, is wanted. Thus the lower and upper endpoints
are kept in their distinct places throughout the mixed transfer.

For later use, subtraction of the exact logarithms in AMT4 gives
\[
\begin{gathered}
 \mathfrak s_h=\alpha+\log(M_\alpha/\vartheta_h),\qquad
 \mathfrak c_h=-\log M_\alpha+3\log\vartheta_h-3\alpha
       +\log\frac{C_h^{\rm tilt}C_\Gamma2^{B_h/2}}{c_hc_\Gamma},\\
 \log(A_k/a_k)=k\mathfrak s_h+(p+1)\log k
                         +B_h\log(k-2)+\mathfrak c_h,\qquad
 E_k^\pm=O_h(k+\log q).
\end{gathered}\tag{AMT22}
\]
The last statement follows from the displayed exact expressions,
since $M,p,B_h$ and the two integral constants are fixed and
$\log D_N=O_h(1+\log(N+1))$. It is a consequence of the original
constants and contains no change of measure.

\subsection{The full original multiplier and its transported observation}
On the existing tensor class $k=4l+1$, retain
\[
 t_j=2j+\tfrac12,\quad\xi_j=c+t_j,\quad
 f_l(S)=\prod_{j=0}^{l-1}(S-c-t_j),\quad
 \beta_k=\frac{(2\pi)^{k/2}}{\sqrt2\Gamma(k/2)},\quad
 dm_{0,k}(y)=\beta_k|f_l(c+iy)|^2d\sigma(y).
 \tag{AMT23}
\]
This is the exact Gamma recurrence identity TWA26 with masses
$(2\pi)^{k/2}$ and $\sqrt{2\pi}$. For $l=0$ both $f_l$ and
$\beta_k$ are one. Write $\mathsf E_\chi:E\to\mathbb C^q$ for
the complete divided Taylor map
$[P]\mapsto(P^{(r)}(\lambda_{a,b})/r!)_{a,b,\,0\le r<e}$.
It is invertible: a vector in its kernel is divisible by the
coprime factors of $\chi$, and dimensions agree. The full
multiplication matrix $U_f$ has, in each primary block, entries
\[
 (U_f)_{r,a}=\begin{cases}
 f_l^{(r-a)}(\lambda)/(r-a)!,&0\le a\le r<e,\\
 0,&a>r.
 \end{cases}\qquad
 \det U_f=\prod_{a,b=0}^k\prod_{j=0}^{l-1}
       [(2a-k)\delta-t_j+i(2b-k)\gamma]^e\ne0.
 \tag{AMT24}
\]
Indeed $k$ is odd, so every $\lambda_{a,b}$ has nonzero imaginary
part, while all $\xi_j$ are real and distinct. The inverse in each
block is the finite Taylor inverse of its nonzero constant term.
The exact polynomial map and its inverse on its stated range are
\[
 \mathcal P_N\longrightarrow f_l\mathcal P_N\subseteq\mathcal P_{N+l},
 \quad P\longmapsto f_lP,
 \qquad Q\longmapsto Q/f_l.
 \tag{AMT25}
\]
It sends the original fibre $J_NP=z$ bijectively to the product
jet fibre $(\mathsf E_fQ,\mathsf E_\chi Q)=(0,U_f\mathsf E_\chi z)$.
Vanishing of the $l$ distinct $f$-values proves divisibility by
$f_l$, and the inverse of $U_f$ proves the converse fibre statement.
The original relation space maps to $f_l\chi\mathcal P_{N-q}$.
On the free $\chi$ jet coordinate $Y$ of this constrained product
fibre, the exact observed map, its kernel, and a section are
\[
\boxed{\begin{gathered}
 \widetilde\Lambda=\Lambda\mathsf E_\chi^{-1}U_f^{-1},\qquad
 \widetilde\Lambda U_f\mathsf E_\chi=\Lambda,\\
 \ker\widetilde\Lambda=U_f\mathsf E_\chi K,\qquad
 \widetilde I=U_f\mathsf E_\chi I,\qquad
 \widetilde S_*=U_f\mathsf E_\chi S_*,\quad
 \widetilde\Lambda\widetilde S_*=I_B.
\end{gathered}}\tag{AMT26}
\]
Both inclusions in the kernel equality follow by applying the
displayed invertible map and its inverse. This retains the full
original amplitude in $\Lambda$ and every derivative in $U_f$.
The resulting kernel is proved by transport; it has not been
assigned the coordinates of $\mathsf E_\chi K$.

For an explicit metric on this transported space, let $b_n(y)$ be
the fixed monic Gamma polynomials and put
$\pi_n(S)=i^n b_n((S-c)/i)$, whose leading coefficient is one
and squared norm is $\gamma_n=\sqrt{2\pi}(2n)!/4^n$.
The evaluation kernels are the full finite sums
\[
 \mathscr K_{\chi,L}=\sum_{n=0}^L
   \frac{\mathsf E_\chi\pi_n(\mathsf E_\chi\pi_n)^*}{\gamma_n},\quad
 \mathscr K_{f,L}=\sum_{n=0}^L
   \frac{\mathsf E_f\pi_n(\mathsf E_f\pi_n)^*}{\gamma_n},\quad
 \mathscr K_{\chi,f;L}=\sum_{n=0}^L
   \frac{\mathsf E_\chi\pi_n(\mathsf E_f\pi_n)^*}{\gamma_n}.
\]
The product evaluation is onto for $L=N+l\ge q+l-1$, because
a polynomial of degree below $q+l$ vanishing at all its divided
jets is divisible by the degree-$q+l$ product $f_l\chi$ and is
zero. Thus its positive kernel has positive Schur complement
\[
\begin{gathered}
 C_N=\mathscr K_{\chi,N+l}-\mathscr K_{\chi,f;N+l}
            \mathscr K_{f,N+l}^{-1}\mathscr K_{f,\chi;N+l}>0,\\
 U_N=(\mathsf E_\chi\pi_{N+j}/\sqrt{\gamma_{N+j}})_{j=1}^l,
 \quad V_N=\mathscr K_{\chi,f;N+l}\mathscr K_{f,N+l}^{-1/2},\quad
 C_N=\mathscr K_{\chi,N}+U_NU_N^*-V_NV_N^*.
\end{gathered}\tag{AMT27}
\]
Empty $f$ blocks have the identity convention. The inverse of the
joint positive kernel, evaluated on $(0,Y)$, gives the minimum
$Y^*C_N^{-1}Y$: block elimination of the $f$-block proves this
identity with every mixed entry present. Combining it with AMT25
and the coefficient $\beta_k$ proves the exact tensor-source metrics
\[
\boxed{\begin{gathered}
 G_N^0=\beta_k\mathsf E_\chi^*U_f^*C_N^{-1}U_f\mathsf E_\chi,\qquad
 H_{K,N}^0=\beta_k\widetilde I^*C_N^{-1}\widetilde I,\\
 Q_N^0=\beta_k(\widetilde\Lambda C_N\widetilde\Lambda^*)^{-1},\qquad
 S_N^0=\mathsf E_\chi^{-1}U_f^{-1}C_N\widetilde\Lambda^*
                   (\widetilde\Lambda C_N\widetilde\Lambda^*)^{-1}.
\end{gathered}}\tag{AMT28}
\]
Direct multiplication verifies the section and its orthogonality.
The formulas hold on the original $E,K,B$, with the transported
observation exactly as in AMT26. No block or mixed product was
dropped from $C_N$.

The same coefficient identity also supplies a quantitative comparison
directly with this original tensor source. Define the positive exact
finite factors
\[
 \Pi_l=\prod_{j=0}^{l-1}t_j^2,\qquad
 U_{l,N}^{\rm poly}=\prod_{j=0}^{l-1}[4(N+j+1)^2+t_j^2],\qquad
 \widetilde L_N=L_N+\log(U_{l,N}^{\rm poly}/\Pi_l).
 \tag{AMT29}
\]
Pointwise $|f_l(c+iy)|^2\ge\Pi_l$. Also
$\|(iy-t)P\|_\sigma^2=\|yP\|_\sigma^2+t^2\|P\|_\sigma^2$,
so the degree-specific multiplication estimate in AMT6, applied
successively, gives
$\|f_lP\|_\sigma^2\le U_{l,N}^{\rm poly}\|P\|_\sigma^2$.
Writing $H_N^0$ for the unchanged $m_{0,k}$ source Gram yields
\[
 \beta_k\Pi_lH_N^\sigma\preceq H_N^0
                  \preceq\beta_k U_{l,N}^{\rm poly}H_N^\sigma,
 \qquad
 \frac{a_k}{D_N\beta_k U_{l,N}^{\rm poly}}H_N^0
       \preceq H_N^{\rm ar}\preceq\frac{A_k}{\beta_k\Pi_l}H_N^0.
 \tag{AMT30}
\]
The second pair follows by using the first upper bound for its lower
inequality and the first lower bound for its upper inequality.
Minimizing and restricting with the same $J_N,I,\Lambda$ gives
AMT30 for all three metrics in AMT28. Its common upper scalar is
$A_k/(\beta_k\Pi_l)$; it is retained before its four signed copies
cancel. The same determinant proof as AMT12--17 therefore gives
\[
\begin{gathered}
 \widetilde E_k^+=E_k^+
   +\log(U_{l,q-1}^{\rm poly}U_{l,q}^{\rm poly}/\Pi_l^2),\qquad
 \widetilde E_k^-=E_k^-
   +\log(U_{l,2q-1}^{\rm poly}U_{l,2q}^{\rm poly}/\Pi_l^2),\\
 \boxed{-r\widetilde E_k^+\le F_K^{\rm ar}-F_K^0
                 \le r\widetilde E_k^-,\qquad
        -b\widetilde E_k^+\le F_B^{\rm ar}-F_B^0
                 \le b\widetilde E_k^-.}
\end{gathered}\tag{AMT31}
\]
For an explicit bound on the extra widths, $t_j\ge1/2$ gives
\[
 0\le\log(U_{l,N}^{\rm poly}/\Pi_l)
   =\sum_{j=0}^{l-1}\log\left(1+\frac{4(N+j+1)^2}{t_j^2}\right)
   \le l\log[1+16(N+l)^2].
 \tag{AMT32}
\]
Thus $\widetilde E_k^\pm=O_h(k\log(q+k))$ on the stated endpoints.
This is a proved polynomial-source estimate for the exact original
$f_l$, with no inference that its multiplication preserves $K$.

\subsection{The evaluated original simple-quartet arithmetic kernel}
Now take the original simple-quartet class $m=1$, $k=4l+1$,
$l\ge1$, $q=(k+1)^2$. The parameter $u$ is precisely the period
parameter in AMT3, not the real integration coordinate $y$.
RPC21 and RPC23 give an explicit finite radius $R_*$ for the
unchanged quartet and arithmetic unit. For reference its exact
constants and cases are recorded here. Set
\[
\begin{gathered}
 \beta=2(\gamma^2-\delta^2),\quad\eta=(\delta^2+\gamma^2)^2,
 \quad a=v_h(\tfrac12+\delta+i\gamma)\ne0,\\
 c_3=\frac{a^{-2}-\overline a^{-2}}{4i\delta\gamma},\qquad
 c_1=\frac{a^{-2}+\overline a^{-2}}2-(\delta^2-\gamma^2)c_3,\\
 g_j=5^{j/5-1}e^{\pi ij/5}\Gamma(j/5)\ (1\le j\le4),\quad
 g_+=\max_j|g_j|,\quad g_-^{-1}=\max_j|g_j|^{-1},\quad
 \kappa=\max_{1\le j<4}|g_{j+1}/g_j|,\\
 E_j^{\rm ray}=\frac85\int_0^\infty r^j(r^3/3+r)
                  e^{-r^5/5+r^3/3+r}\,dr\ (0\le j\le3),\quad
 C_T=2\left(\sum_{j=0}^3(E_j^{\rm ray})^2\right)^{1/2},\\
 C_K=4C_Tg_-^{-1}+2g_+g_-^{-1}+2g_+g_-^{-2}C_T,\quad
 C_{K,3}=C_K(3\kappa^2+3\kappa C_K+C_K^2),\\
 \epsilon_*=\min\{1,(2C_Tg_-^{-1})^{-1}\},\qquad B_*=\beta+\eta.
\end{gathered}\tag{AMT33}
\]
Here $g_-^{-2}=(g_-^{-1})^2$. All integrals converge by their
negative fifth-degree exponent, and $(c_1,c_3)\ne(0,0)$ by the
original unit values. The exact radius is
\[
 R_*=
 \begin{cases}
 \displaystyle\left[\max\left\{1,\frac{B_*}{\epsilon_*},
 \frac{2(|c_3|C_{K,3}B_*+|c_1|(\kappa+C_K))}{|c_3g_4/g_1|}\right\}\right]^{5/2},
                  &c_3\ne0,\\[6pt]
 \displaystyle\left[\max\left\{1,\frac{B_*}{\epsilon_*},
                     \frac{2C_KB_*}{|g_2/g_1|}\right\}\right]^{5/2},
                  &c_3=0.
 \end{cases}\tag{AMT34}
\]
The complete period calculation RPC16--25 proves, for every original
branch and every $|u|\ge R_*$, a nonzero original commutator entry:
its $(4,1)$ entry in the first case and its $(2,1)$ entry in the
second. RPC14--15 identifies those entries with the exact test for
the original rank-four quadric's cyclic invariance. Thus its actual
projective orbit has five elements throughout this domain. OQX1--3,
using the product of precisely these five original quadrics, proves
\[
 \boxed{r=8k-16,\qquad b=k^2-6k+17,
        \qquad m=1,\ k=4l+1,\ l\ge1,\ |u|\ge R_*.}
 \tag{AMT35}
\]
This is the evaluated original period outcome; no value of the
amplitude or of a period coefficient is chosen. The radius depends
on the fixed one-factor quartet and unit and is independent of $k$.

For the same original packet, OKV22 proves for both existing source
orders $s=1,k$, and for the actual $K$ of any rank,
\[
\begin{gathered}
 D_*=\max\{3,\sqrt{\delta^2+\gamma^2}\},\quad
 C(\delta,\gamma)=12+4\log(2D_*),\\
 Z_q=1+2(q+1)q^2 2^{6q}(2D_*q)^{4q},\qquad
 0\le F_K^{(s)}\le2r\log Z_q
                \le2r C(\delta,\gamma)q\log(q+1),
 \quad F_K^{(1)}=F_K^\sigma,\quad F_K^{(k)}=F_K^0.
\end{gathered}\tag{AMT36}
\]
Its proof includes the exact source coefficient map and the unchanged
kernel inclusion. Combining it with the already proved AMT17--20
gives the finite arithmetic conclusion
\[
\boxed{\begin{split}
 \max\{0,F_K^\sigma-C_K^-\}
       &\le F_K^{\rm ar}\le F_K^\sigma+C_K^+,\\
 0\le F_K^{\rm ar}
       &\le U_K:=2r\log Z_q+\widehat C_K^+\\
       &\le2r C(\delta,\gamma)q\log(q+1)+rE_k^-.
\end{split}}\tag{AMT37}
\]
The cap symbols $C_K^\pm,\widehat C_K^\pm$ in this formula mean
the endpoint sums AMT16; the constant $C_K$ without a superscript
in AMT33 belongs only to the explicit period radius. They have
different stated indices and are not identified.
The alternative finite tensor-source bound
$F_K^{\rm ar}\le2r\log Z_q+r\widetilde E_k^-$ also follows from
AMT31. Formula AMT37 retains the sharper fixed-source widths.

On AMT35 the exact ranks give
\[
 \boxed{\quad F_K^{\rm ar}-F_K^\sigma=O_h(k^2),\qquad
 F_K^{\rm ar}=O_h(kq\log(q+1)),\qquad
 \frac{F_K^{\rm ar}}{q^2}\longrightarrow0.\quad}
 \tag{AMT38}
\]
For the first estimate use $r=8k-16$ in AMT20 and AMT22,
with $\log q=2\log(k+1)$. The second follows from AMT37.
For the limit, divide its last line by $q^2$: its first term is
at most $16C(\delta,\gamma)k\log(q+1)/q\to0$, and its second
is $O_h(k^2/q^2)\to0$. Nonnegativity gives the limit by the two
bounds. These estimates are uniform over the entire original
domain $|u|\ge R_*$, since none of their right sides depends on $u$.
The tensor comparison likewise gives
$F_K^{\rm ar}-F_K^0=O_h(k^2\log q)=o(q^2)$ from AMT31--32.

\subsection{The original boundary receives the full baseline coefficient}
Before using a limit, the exact finite arithmetic identity and its
consequence are
\[
\boxed{\begin{gathered}
 F_B^{\rm ar}=\mathcal B_k^\sigma+\delta_k^\sigma-F_K^{\rm ar},\\
 \max\{0,\mathcal B_k^\sigma+\delta_k^\sigma-U_K\}
    \le F_B^{\rm ar}\le\mathcal B_k^\sigma+\delta_k^\sigma,\\
 \max\{0,\mathcal B_k^\sigma+\underline\delta_k^\sigma-U_K\}
    \le F_B^{\rm ar}\le\mathcal B_k^\sigma+\overline\delta_k^\sigma.
\end{gathered}}\tag{AMT39}
\]
The first identity follows from AMT10 and AMT15; AMT37 and
$F_K^{\rm ar},F_B^{\rm ar}\ge0$ prove the first interval, and
AMT21 proves the second. The correlated boundary return also has
the sharper exact interval AMT18, with every original signed
total correction present.

BSL13--19 supplies the complete original finite baseline and its
proved coefficient
\[
 \mathcal B_k^\sigma=V_q+\eta_q+\eta_{q+1}-\eta_1,\qquad
 \frac{\mathcal B_k^{\rm ar}}{q^2}\longrightarrow
 C_B^{\rm base}:=9-8\log2+\mathcal F(2,\pi)
                      >\frac{769}{17010}>0.
 \tag{AMT40}
\]
Here $V_q$, the three distinct original relation errors $\eta_r$,
and the energy $\mathcal F$ are exactly the BSL/HBL/GEL quantities,
whose complete definitions and proofs are in the retained closure.
The superscript on $C_B^{\rm base}$ only distinguishes that existing
baseline coefficient from the endpoint caps $C_B^\pm$ in AMT16.
Its finite BSL14 enclosure can be substituted directly for
$\mathcal B_k^\sigma$ in AMT39, keeping the signs of all three
$\eta$ terms; no new evaluation is needed for that substitution.
Subtract AMT38 from the exact original identity
$F_B^{\rm ar}=\mathcal B_k^{\rm ar}-F_K^{\rm ar}$. It proves
\[
 \boxed{\quad
 \frac{F_B^{\rm ar}}{q^2}\longrightarrow C_B^{\rm base}>0,
 \qquad
 F_K^{\rm ar}+F_B^{\rm ar}=\mathcal B_k^{\rm ar},
 \quad m=1,\ k=4l+1,\ l\to\infty,\ |u|\ge R_*.
 \quad}\tag{AMT41}
\]
The convergence is uniform over that original period domain by
AMT38 and because the full-source baseline is independent of $u$.
Thus the proved original kernel estimate now passes to the actual
arithmetic kernel, and the original arithmetic boundary carries the
entire positive $q^2$ coefficient. All earlier full-sum identities
remain exact, with their individual mixed contributions now receiving
AMT37--41 and their signed finite corrections AMT17--21.

\subsection{The same comparison evaluates the mixed Gamma return on the $kq$ scale}
The first inequality pair AMT30 directly compares the two existing
Gamma sources on the same original coefficient space. Apply its
proved minimum and restriction maps AMT7--9 to $X=E,K,B$, writing
$H_{E,N}^a=G_N^a$, $H_{B,N}^a=Q_N^a$, and
$r_E=q,r_K=r,r_B=b$. At every original endpoint it gives
\[
 \beta_k\Pi_l H_{X,N}^\sigma\preceq H_{X,N}^0
       \preceq\beta_k U_{l,N}^{\rm poly}H_{X,N}^\sigma.
 \tag{AMT42}
\]
Define the exact relative determinant and the positive degree width
\[
\begin{gathered}
 z_{X,N}=\log\frac{\det H_{X,N}^0}{\det H_{X,N}^\sigma}
                              -r_X\log(\beta_k\Pi_l),\qquad
 \omega_N^\Gamma=\log(U_{l,N}^{\rm poly}/\Pi_l),\\
 0\le z_{X,N}\le r_X\omega_N^\Gamma,\qquad
 z_{K,N}+z_{B,N}=z_{E,N},\\
 \Omega_k^{\rm low}=\omega_{q-1}^\Gamma+\omega_q^\Gamma,\qquad
 \Omega_k^{\rm high}=\omega_{2q-1}^\Gamma+\omega_{2q}^\Gamma.
\end{gathered}\tag{AMT43}
\]
Congruence and multiplication of the eigenvalues in AMT42 prove
the two inequalities. The determinant identity AMT10, applied to
the tensor metric AMT28 as well, proves the common sum. These
facts preserve the complete original kernel and boundary rather
than assigning separate arbitrary errors to their determinants.

Use the precise return orientation
\[
 \Delta_X^\Gamma=F_X^\sigma-F_X^0\quad(X=K,B),\qquad
 \mathcal H_k=\mathcal B_k^\sigma-\mathcal B_k^0=-\mathfrak T_k.
\]
The last equality is the original TWA30/BSL18 sign. Substituting
all four $z$ terms, and canceling $\log(\beta_k\Pi_l)$ only
after the signed sum, proves
\[
\boxed{\begin{split}
 \Delta_X^\Gamma
   &=-z_{X,q-1}-z_{X,q}+z_{X,2q-1}+z_{X,2q},\\
 -r_X\Omega_k^{\rm low}
   &\le\Delta_X^\Gamma\le r_X\Omega_k^{\rm high},\\
 \Delta_K^\Gamma+\Delta_B^\Gamma&=\mathcal H_k,\\
 \mathcal H_k-r\Omega_k^{\rm high}
   &\le\Delta_B^\Gamma\le\mathcal H_k+r\Omega_k^{\rm low}.
\end{split}}\tag{AMT44}
\]
AMT29 displays every factor of both $\Omega$ values explicitly;
AMT32 gives the finite bounds
\[
\begin{split}
 0\le\Omega_k^{\rm low}
 &\le l\log[1+16(q-1+l)^2]+l\log[1+16(q+l)^2],\\
 0\le\Omega_k^{\rm high}
 &\le l\log[1+16(2q-1+l)^2]+l\log[1+16(2q+l)^2].
\end{split}\tag{AMT45}
\]
In particular both are $O(k\log(q+1))$ on the original
$m=1$ class. Its actual rank $r=8k-16$ therefore proves directly
\[
 \boxed{\Delta_K^\Gamma=O_h(k^2\log(q+1))=o(kq),\qquad
 \frac{F_K^{\rm ar}-F_K^\sigma}{kq}\longrightarrow0,\qquad
 \frac{F_K^{\rm ar}-F_K^0}{kq}\longrightarrow0.}
 \tag{AMT46}
\]
For the first equality divide the AMT44 bounds by $kq$ and use
$q=(k+1)^2$, so $k\log(q+1)/q\to0$. The second limit uses
AMT38, since $k^2/(kq)=k/q\to0$. The last uses the exact
identity $F_K^{\rm ar}-F_K^0=\Delta_K^\sigma+\Delta_K^\Gamma$.
All bounds are uniform for $|u|\ge R_*$.

For completeness the entire finite arithmetic boundary return is
also evaluated by those same quantities:
\[
\boxed{\begin{split}
 F_B^{\rm ar}-F_B^0
   &=\mathcal H_k+\delta_k^\sigma
                         -(F_K^{\rm ar}-F_K^0),\\
 \mathcal H_k+\delta_k^\sigma-r\widetilde E_k^-
   &\le F_B^{\rm ar}-F_B^0
             \le\mathcal H_k+\delta_k^\sigma+r\widetilde E_k^+.
\end{split}}\tag{AMT47}
\]
This uses AMT31 with its exact endpoint constants
$\widetilde E_k^+=E_k^++\Omega_k^{\rm low}$ and
$\widetilde E_k^-=E_k^-+\Omega_k^{\rm high}$, rather than a
new source comparison. Both tensor and arithmetic source masses
and the full original resultant have already been retained in
AMT23--30 before these signed cancellations.

The proved GEL18--19 constant is
\[
 \mathcal C_\Gamma=2\int_{a_*^2}^{b_*^2}\log x\,
                       \rho_{2,\pi}(x)\,dx-4(2\log2-1)>0,
 \qquad \frac{\mathfrak T_k}{lq}\longrightarrow\mathcal C_\Gamma.
 \tag{AMT48}
\]
Here $a_*,b_*$ and $\rho_{2,\pi}$ are exactly the proved GEL8/EIQ
equilibrium endpoints and density; their complete definitions and
integral proof remain in that provider. The endpoint letters here
avoid the already fixed original period parameter $u$. GEL18a
gives $4\log(27/(8\pi))\le\mathcal C_\Gamma\le6$ and proves
the strict positivity. TWA15 also proves the actual source return
$\delta_k^\sigma/(kq)\to0$. As $l/k\to1/4$, AMT44, AMT46
and AMT47 now prove
\[
\boxed{\begin{gathered}
 \frac{F_B^\sigma-F_B^0}{kq}\longrightarrow-\frac{\mathcal C_\Gamma}{4},
 \qquad
 \frac{F_B^{\rm ar}-F_B^0}{kq}\longrightarrow-\frac{\mathcal C_\Gamma}{4},
 \qquad
 \frac{F_B^{\rm ar}-F_B^\sigma}{kq}\longrightarrow0,\\
 \frac{F_B^\sigma-F_B^0}{\mathcal H_k}\longrightarrow1,
 \qquad
 \frac{F_B^{\rm ar}-F_B^0}{\mathcal H_k}\longrightarrow1,
 \quad m=1,\ k=4l+1,\ l\to\infty,\ |u|\ge R_*.
\end{gathered}}\tag{AMT49}
\]
Indeed the first two limits subtract an $o(kq)$ kernel return
from $\mathcal H_k=-\mathfrak T_k$, whose quotient by $kq$
tends to $-\mathcal C_\Gamma/4$. The third follows from
$\Delta_B^\sigma=\delta_k^\sigma-\Delta_K^\sigma$.
The nonzero limit of $\mathcal H_k/(kq)$ makes its denominator
nonzero for all sufficiently large original $k$ and proves the
last two ratios by division. The uniformity in $u$ follows from
the uniform kernel bounds and the period-independent total returns.
Thus the full existing source comparison evaluates the mixed
Gamma return also at its $kq$ scale, with its finite signed
remainders AMT44 and AMT47 still present.

\end{document}
